% Level 1, Week 1 -- Robots Need Instructions (Scratch)
\lesson{1}{Robots Need Instructions}{A program is a list of exact steps, done in order.}
\lessoninfo{Concepts: sequence (I), algorithm (I), sprite / stage / green flag (I)}{Tools: Scratch 3 offline app}{CSTA: 1B-AP-08, 1B-AP-10, 1B-NI-05}

\begin{objectives}
  \item The student can write an algorithm as exact steps in order so that a \enquote{robot} draws the right picture. (L1.G1) [L1.W01.1]
  \item The student can point to the stage, a sprite, the blocks, the code area and the green flag in the Scratch app, and run a script. (L1.G1) [L1.W01.2]
  \item The student can find and fix a script whose steps are in the wrong order. (L1.G6) [L1.W01.3]
  \item The student can explain our class agreements and why we use a nickname and keep projects on our own Mac. (L1.G7, L1.G8) [L1.W01.4]
\end{objectives}

\begin{materials}
  \item \textbf{One-time setup (the week before, about 20 minutes):} install the Scratch app on every Mac (see the teacher note \enquote{Installing Scratch} below) and make a folder called \texttt{Scratch Projects} on each Mac's Desktop.
  \item Print: workbook chapter 1 and homework 1 (one per student). The class agreements are on the first page of the chapter.
  \item For the Parent-robot game: squared (graph) paper, one thick marker, a pad of sticky notes per pair, and the three picture cards in this plan (cut them apart, or just point at them).
  \item A timer (phone or kitchen timer) for driver/navigator swaps.
  \item Optional: the Interest and Comfort Survey in the Assessments booklet (not graded, about 10 minutes; it can go home with the homework).
\end{materials}

\begin{teachernote}[Installing Scratch (do this before the first session)]
Scratch is free and works with no account and no internet once it is installed.
\begin{enumerate}[leftmargin=1.4em,itemsep=1pt]
  \item On the Mac, open the \menu{App Store}, type \texttt{Scratch} in the search box and pick the app called \enquote{Scratch 3} (the publisher is the Scratch Foundation, price: free). Click \menu{Get}, then \menu{Install}. If you prefer not to use the App Store, open \texttt{scratch.mit.edu/download} in Safari, choose \menu{macOS} and follow the steps there.
  \item Open the app once: press \keys{cmd}+\keys{Space}, type \texttt{Scratch}, press \keys{Return}. If it asks whether to share usage data, choose the option that does not share.
  \item Right-click the Scratch icon in the Dock and choose \menu{Options} $\rightarrow$ \menu{Keep in Dock}, so students can find it.
  \item Check: you see a cat in the white stage (top right) and a green flag above it. Click the flag: nothing happens yet (no code), and that is correct.
  \item Make the save folder: click the Desktop, press \keys{shift}+\keys{cmd}+\keys{N}, name it \texttt{Scratch Projects}.
\end{enumerate}
We never use the \enquote{Share} or online parts of Scratch in this course. The offline app has no sign-in at all.
\end{teachernote}

\needspace{16\baselineskip}
\section{At a glance}
\begin{timeplan}
  \timeslot{0--8}{Warm-up}{Welcome and class agreements (3 min). Parent-robot drawing, round 1 (5 min).}
  \timeslot{8--18}{Concept}{Debrief the robot: what is an algorithm? Steps must be exact and in order. Round 2 with a swapped sticky note (a bug!). Nicknames and offline safety.}
  \timeslot{18--28}{Guided demo}{Open Scratch together. Tour: stage, sprite, blocks, code area, green flag. Build \enquote{Hello, Cat!} on the big screen.}
  \timeslot{28--52}{Hands-on build}{Pairs: driver/navigator, swap every 8 minutes. Exercise 1.1, then the Mixed-Up Jump bug hunt (1.2), then My Sprite Show (1.3). Last 5 minutes free-build. Save to \texttt{Scratch Projects}.}
  \timeslot{52--60}{Wrap-up}{Show-and-tell (1 minute each), exit ticket, reflection, hand out homework 1.}
\end{timeplan}

\section{Warm-up (0--8 min)}
\begin{teacherscript}[Welcome and agreements]
\enquote{Welcome to coding club! Over the next weeks you will make games and art on the computer. Before we start, let's agree how we work together.} Read the six agreements on the first page of the workbook chapter. After each one, ask: \enquote{Thumbs up if you agree.} Everyone (you too) writes their nickname at the bottom. \enquote{Coders use nicknames. Today, pick a coding nickname you like. It can be funny, like Pixel Panda.}
\end{teacherscript}

\begin{unplugged}[Parent-robot drawing]
\textbf{Setup:} you are the robot. Put squared paper on the table and a dot near the bottom-left. The robot starts on the dot, facing the top of the page. The robot only understands five commands (write them on a sticky note so everyone sees them):
\begin{center}
\texttt{PEN DOWN} \quad \texttt{PEN UP} \quad \texttt{FORWARD n} (squares) \quad \texttt{TURN RIGHT} \quad \texttt{TURN LEFT}
\end{center}
(Turns are always a quarter turn, on the spot.)

\textbf{Play:} show picture card A only to the students (hide it from the robot). In pairs, they write the steps, \textbf{one command per sticky note}, and stick the notes in a column. Then they read the notes aloud one at a time and you draw exactly what they say. Be a very literal robot: if they forget \texttt{PEN DOWN}, move your marker in the air and draw nothing. If they say \enquote{draw a square}, answer in a robot voice: \enquote{DOES NOT COMPUTE.} Compare your drawing to the card. Faster groups do card B, then card C.
\end{unplugged}

\begin{center}
\begin{tikzpicture}[scale=0.5]
  % Card A: square
  \draw[kcgray!50] (0,0) grid (5,5);
  \draw[kcblue,line width=2pt] (1,1) rectangle (4,4);
  \fill[kcvermillion] (1,1) circle (0.22);
  \draw[-{Stealth},kcvermillion,line width=1.2pt] (1,1) -- (1,1.9);
  \node[below] at (2.5,0) {\small Card A: square (3)};
  % Card B: stairs
  \begin{scope}[xshift=9cm]
  \draw[kcgray!50] (0,0) grid (5,5);
  \draw[kcgreen,line width=2pt] (1,1) -- (1,2) -- (2,2) -- (2,3) -- (3,3) -- (3,4) -- (4,4);
  \fill[kcvermillion] (1,1) circle (0.22);
  \draw[-{Stealth},kcvermillion,line width=1.2pt] (1,1) -- (1,1.9);
  \node[below] at (2.5,0) {\small Card B: stairs};
  \end{scope}
  % Card C: letter T
  \begin{scope}[xshift=18cm]
  \draw[kcgray!50] (0,0) grid (5,5);
  \draw[kcpurple!80!black,line width=2pt] (2,1) -- (2,4) (1,4) -- (3,4);
  \fill[kcvermillion] (2,1) circle (0.22);
  \draw[-{Stealth},kcvermillion,line width=1.2pt] (2,1) -- (2,1.9);
  \node[below] at (2.5,0) {\small Card C: letter T (tricky!)};
  \end{scope}
\end{tikzpicture}

{\small The dot is the start; the arrow shows which way the robot faces at the start.}
\end{center}

\section{Concept (8--18 min)}
\begin{teacherscript}
\enquote{What went wrong when the robot drew? (Let them answer.) The robot did exactly what you said, not what you meant. A list of exact steps in the right order is called an \textbf{algorithm}. When a computer follows an algorithm, we call it a \textbf{program}. Computers are like our robot: super fast, but they never guess.}

Now play round 2 with card A. While nobody is looking, swap two sticky notes (for example a \texttt{TURN RIGHT} and a \texttt{FORWARD 3}). Draw it. \enquote{Uh-oh, the robot drew something strange. Same steps, but two are swapped. A mistake in a program is called a \textbf{bug}, and finding and fixing it is called \textbf{debugging}. Bugs are normal. Every coder makes them every day. Can you debug my sticky notes?} Let them fix the order and redraw.
\end{teacherscript}

\begin{analogy}[A recipe for pancakes]
\enquote{An algorithm is like a recipe. If the recipe says \enquote{flip the pancake} before \enquote{pour the batter in the pan}, you get a mess, even though every step is a real step. The order matters as much as the steps.}
\end{analogy}

\begin{safety}[Nicknames and offline projects]
\enquote{In our projects we use our coding nickname, never our real full name, home address, school name, birthday or photos of ourselves. The Scratch app on this Mac works without the internet and without an account. Our projects stay on this computer. If any app or website ever asks for your real name, address or a password, stop and get a grown-up.} Ask each student to say their nickname once; that is the name they type in projects.
\end{safety}

\section{Guided demo (18--28 min)}
\begin{teacherscript}
Open Scratch on your Mac (click the cat in the Dock) and let the students sit where they can see. \enquote{This is Scratch. Let me show you the five things you need today.} Point to each part of the screen as you say it (the workbook has the same labeled picture):
\begin{itemize}[leftmargin=1.2em,itemsep=0pt]
  \item \textbf{Stage} (white area, top right): where things happen. \textbf{Sprite}: a character on the stage; the cat is a sprite.
  \item \textbf{Block categories} (colored circles, far left) and the \textbf{blocks} next to them.
  \item \textbf{Code area} (big middle space): where we snap blocks together into a \textbf{script}.
  \item \textbf{Green flag} (above the stage) starts the program; the \textbf{red stop sign} stops it.
\end{itemize}
\end{teacherscript}

\begin{workedexample}[Hello, Cat! (Exercise 1.1)]
Narrate each click out loud. Build this script, slowly:
\par\noindent\begin{minipage}[t]{0.56\linewidth}
\begin{enumerate}[leftmargin=1.4em,itemsep=1pt]
  \item Click the yellow \menu{Events} circle. Drag \enquote{when green flag clicked} into the code area.
  \item Click \menu{Motion}. Drag \enquote{go to x: 0 y: 0} and snap it under the hat block (you will see a gray shadow when it is about to snap). Click the numbers and type \texttt{0} and \texttt{0}: that is the center of the stage.
  \item Click \menu{Looks}. Snap \enquote{say Hello! for 2 seconds} underneath. Click the word \texttt{Hello!} and type \texttt{Hi! I am Pixel Panda} (use your nickname).
  \item From \menu{Motion}, snap \enquote{move 10 steps} and change 10 to \texttt{100}.
  \item From \menu{Sound}, snap \enquote{play sound Meow until done}.
  \item From \menu{Looks}, snap \enquote{say Hello! for 2 seconds} and type \texttt{Bye!}.
  \item Click the green flag. Ask: \enquote{Which step happened first? Which came last?}
\end{enumerate}
\end{minipage}\hfill
\begin{minipage}[t]{0.40\linewidth}\vspace{4pt}\centering
\begin{scratch}[scale=0.68]
\blockinit{when \greenflag clicked}
\blockmove{go to x: \ovalnum{0} y: \ovalnum{0}}
\blocklook{say \ovalnum{Hi! I am Pixel Panda} for \ovalnum{2} seconds}
\blockmove{move \ovalnum{100} steps}
\blocksound{play sound \selectmenu{Meow} until done}
\blocklook{say \ovalnum{Bye!} for \ovalnum{2} seconds}
\end{scratch}
\end{minipage}

Then show how to save: \menu{File} $\rightarrow$ \menu{Save to your computer}, choose the \texttt{Scratch Projects} folder on the Desktop, type the name \texttt{hello-cat} and click \menu{Save}. To remove a wrong block: drag it back onto the block area on the left and let go.
\end{workedexample}

\section{Hands-on build (28--52 min)}
\begin{pairroles}[Driver and navigator]
\enquote{Coders often work in pairs. The \textbf{driver} holds the mouse and does the clicking. The \textbf{navigator} reads the workbook steps aloud, watches for bugs and says what to do next. Navigators do not grab the mouse: they point and explain. When the timer rings, swap seats.} Set the timer for 8 minutes.

With 1 student: you are the navigator for the first 8 minutes, then swap and let your child tell you what to click. With 2: one pair. With 3: a rotating trio (driver, navigator, tester who clicks the green flag and checks the workbook checklist); rotate at each timer. With 4: two pairs.
\end{pairroles}

Pairs work through the workbook: \textbf{Exercise 1.1} Hello, Cat! (copy of the demo, using their own nickname), \textbf{Exercise 1.2} Mixed-Up Jump (bug hunt below), \textbf{Exercise 1.3} My Sprite Show (choose a sprite and a backdrop, write a 5-step show). Walk round and ask \enquote{What will happen first? Click the flag and see.}

\begin{bughunt}[Mixed-Up Jump: steps in the wrong order]
In Exercise 1.2 the cat should say \enquote{Watch me jump!}, jump up, land back on the ground and say \enquote{Ta-da!}. The workbook shows the script with two blocks in the wrong order: the cat goes \textbf{down} first (into the ground) and then up. Students build it exactly as printed, click the flag, describe what they see, then drag the two blocks into the right order.
\textbf{Debugging strategy to teach:} 1. Say what you expected. 2. Say what actually happened. 3. Read the script one block at a time, pointing with your finger, and act it out (\enquote{change y by $-50$: down!}). 4. Change one thing and test again.
\end{bughunt}

\begin{freebuild}[My Sprite Show]
Last 5 minutes: \enquote{Add anything you like to your show: a new backdrop, another sound, a different sprite, a longer story. The only rule: click the green flag often to test.} Remind them to save (\menu{File} $\rightarrow$ \menu{Save to your computer}).
\end{freebuild}

\section{Wrap-up (52--60 min)}
\begin{teacherscript}
\enquote{Let's have a mini show. Each pair clicks the green flag on one project, and the rest of us watch.} After each show, ask the audience: \enquote{What was the first step? What was the last?}
\end{teacherscript}

\begin{formativecheck}[Exit ticket]
\textbf{Look for:} every student can start a script with the green flag, snap blocks under a hat block, and explain a fix in Exercise 1.2.

\textbf{Ask each student (out loud or on a sticky note):} \enquote{A robot should put on socks and shoes. Its program is: 1. Put on shoes. 2. Put on socks. What is the bug, and how do you fix it?}
\textbf{Answer:} the steps are in the wrong order; swap them so socks come before shoes.
\end{formativecheck}

\begin{reflection}
Ask: \enquote{What surprised you about giving instructions to a robot?} and \enquote{What is one thing you want to make in Scratch?} Students write one sentence in the reflection box of their workbook. Hand out homework 1 (Walk and Talk) and point out the Core task.
\end{reflection}

\section{Notes for the teacher}
\begin{commonmistakes}
  \item \textbf{Blocks not snapped.} A block that is a little bit apart does not run. Fix: drag it until the gray shadow appears, then let go. Clicking the flag only runs blocks attached under a hat block.
  \item \textbf{Two separate scripts both starting with the flag.} Both run at the same time and things happen at once. Fix: one stack per idea; join them into one column.
  \item \textbf{Typing a real name in the speech bubble.} Gently remind: nickname only.
  \item \textbf{\enquote{It did nothing!}} Often the wrong sprite is selected (the code belongs to the sprite highlighted in the sprite pane). Fix: click the right sprite.
  \item \textbf{\enquote{What does $-100$ mean?}} Negative numbers are new for many 10-year-olds. Say: \enquote{0 is the middle. The minus sign means the other way: y: $-100$ is 100 steps below the middle.} The workbook tip \enquote{Where is x and y?} says the same.
  \item \textbf{The cat wanders off the stage.} Scratch remembers where a sprite ended. That is why we start with \enquote{go to x: 0 y: 0}.
\end{commonmistakes}

\begin{differentiation}
\supportpath Do Exercise 1.1 together with you as navigator and skip 1.3. In the Parent-robot game, use only card A. Let the student drag blocks while you read the words on them.
\extensionpath In My Sprite Show, add a second sprite that answers the first one: use \enquote{wait 2 seconds} so the two speech bubbles take turns, like a conversation.
\end{differentiation}

\begin{stretchchallenge}[A robot dance]
Make the cat do a 6-step dance using only Motion and Looks blocks (\enquote{turn 15 degrees}, \enquote{next costume}, \enquote{change y by 20}, \enquote{wait 0.5 seconds}). Then swap two steps and ask your partner to find the bug.
\end{stretchchallenge}

\begin{teachernote}[If you are short on time]
Do only card A in the Parent-robot game, skip round 2 (instead, point out the Mixed-Up Jump bug later), and leave Exercise 1.3 for free time next week. Do not skip the agreements, the nickname talk or saving: they set up every later session.
\end{teachernote}
